\chapter{2015年浙江卷（理）}

\section{单选题}

\begin{problem}
已知集合 \(P = \{x \mid x^2 - 2x \geqslant 0\}\)，\(Q = \{x \mid 1 < x \leqslant 2\}\)，则 \((\complement_{\R} P) \cap Q =\)（\quad）

\choices
    {\([0,1)\)}
    {\((0,2]\)}
    {\((1,2)\)}
    {\([1,2]\)}
\end{problem}

\begin{answer}
C
\end{answer}

\begin{solution}
由 \(x^2-2x=x(x-2)\geqslant0\)，得 \(P=(-\infty,0]\cup[2,+\infty)\)，所以
\(\complement_{\R}P=(0,2)\). 又 \(Q=(1,2]\)，故
\((\complement_{\R}P)\cap Q=(1,2)\). 因此选 C.
\end{solution}

\begin{problem}
某几何体的三视图如图所示（单位：\(\mathrm{cm}\)），则该几何体的体积是（\quad）

\bitmapfigure[max width=0.55\linewidth]{img/2015/zhejiang_science/q2_fig1.png}

\choices
    {\(8\mathrm{~cm}^{3}\)}
    {\(12\mathrm{~cm}^{3}\)}
    {\(\frac{32}{3}\mathrm{~cm}^{3}\)}
    {\(\frac{40}{3}\mathrm{~cm}^{3}\)}
\end{problem}

\begin{answer}
C
\end{answer}

\begin{solution}
由三视图可知，该几何体由一个底面边长为 \(2\)，高为 \(2\) 的正方体和一个底面边长为 \(2\)，高为 \(2\) 的四棱锥组成. 因此其体积为
\(V=2^3+\dfrac{1}{3}\times2^2\times2=\dfrac{32}{3}\mathrm{~cm}^3\). 所以选 C.
\end{solution}

\begin{problem}
已知 \(\{a_n\}\) 是等差数列，公差 \(d\) 不为零，前 \(n\) 项和是 \(S_n\)，若 \(a_3\)，\(a_4\)，\(a_8\) 成等比数列，则（\quad）

\choices
    {\(a_{1}d > 0\)，\(dS_{4} > 0\)}
    {\(a_{1}d < 0\)，\(dS_{4} < 0\)}
    {\(a_{1}d > 0\)，\(dS_{4} < 0\)}
    {\(a_{1}d < 0\)，\(dS_{4} > 0\)}
\end{problem}

\begin{answer}
B
\end{answer}

\begin{solution}
设 \(a_1=a\). 由 \(a_3\)，\(a_4\)，\(a_8\) 成等比数列，得
\((a+3d)^2=(a+2d)(a+7d)\). 展开并整理，得
\(d(3a+5d)=0\). 由于 \(d\neq0\)，所以 \(a_1=-\dfrac{5}{3}d\)，从而
\(a_1d=-\dfrac{5}{3}d^2<0\).

又 \(S_4=\dfrac{4}{2}(2a_1+3d)=-\dfrac{2}{3}d\)，所以
\(dS_4=-\dfrac{2}{3}d^2<0\). 因此选 B.
\end{solution}

\begin{problem}
命题“\(\forall n \in \N^*\)，\(f(n) \in \N^*\) 且 \(f(n) \leqslant n\)”的否定形式是（\quad）

\choices
    {\(\forall n\in \N^*\)，\(f(n)\notin \N^*\) 且 \(f(n) > n\)}
    {\(\forall n\in \N^*\)，\(f(n)\notin \N^*\) 或 \(f(n) > n\)}
    {\(\exists n_0\in \N^*\)，\(f(n_0)\notin \N^*\) 且 \(f(n_0) > n_0\)}
    {\(\exists n_0\in \N^*\)，\(f(n_0)\notin \N^*\) 或 \(f(n_0) > n_0\)}
\end{problem}

\begin{answer}
D
\end{answer}

\begin{solution}
令 \(A(n)\) 表示命题 \(f(n)\in\N^*\)，令 \(B(n)\) 表示命题 \(f(n)\leqslant n\). 原命题是“对任意 \(n\in\N^*\)，\(A(n)\land B(n)\)”.

否定全称命题得存在 \(n_0\in\N^*\)，使得 \(\neg A(n_0)\lor\neg B(n_0)\). 由于
\(\neg A(n_0)\) 为 \(f(n_0)\notin\N^*\)，\(\neg B(n_0)\) 为 \(f(n_0)>n_0\)，故否定形式为
\(\exists n_0\in\N^*\)，\(f(n_0)\notin\N^*\) 或 \(f(n_0)>n_0\). 因此选 D.
\end{solution}

\begin{problem}
如图，设抛物线 \(y^{2} = 4x\) 的焦点为 \(F\)，不经过焦点的直线上有三个不同的点 \(A\)，\(B\)，\(C\)，其中点 \(A\)，\(B\) 在抛物线上，点 \(C\) 在 \(y\) 轴上，则 \(\triangle BCF\) 与 \(\triangle ACF\) 的面积之比是（\quad）

\bitmapfigure[max width=0.42\linewidth]{img/2015/zhejiang_science/q5_fig1.png}

\choices
    {\(\dfrac{|BF| - 1}{|AF| - 1}\)}
    {\(\dfrac{|BF|^2 - 1}{|AF|^2 - 1}\)}
    {\(\dfrac{|BF| + 1}{|AF| + 1}\)}
    {\(\dfrac{|BF|^2 + 1}{|AF|^2 + 1}\)}
\end{problem}

\begin{answer}
A
\end{answer}

\begin{solution}
设抛物线上点 \(P(x,y)\) 到焦点 \(F(1,0)\) 的距离为 \(PF\). 由 \(y^2=4x\)，有
\(PF^2=(x-1)^2+y^2=(x-1)^2+4x=(x+1)^2\). 因为抛物线上点满足 \(x\geqslant0\)，所以 \(PF=x+1\)，即
\(x=PF-1\).

三点 \(A\)，\(B\)，\(C\) 共线，且 \(C\) 在 \(y\) 轴上. 三角形 \(BCF\) 与 \(ACF\) 以同一直线上的 \(BC\)，\(AC\) 为底边，它们的高相同，因此
\(\dfrac{S_{\triangle BCF}}{S_{\triangle ACF}}=\dfrac{BC}{AC}=\dfrac{x_B}{x_A}=\dfrac{|BF|-1}{|AF|-1}\). 因此选 A.
\end{solution}

\begin{problem}
设 \(A\)，\(B\) 是有限集，定义： \(d(A, B) = \operatorname{card}(A \cup B) - \operatorname{card}(A \cap B)\)，其中 \(\operatorname{card}(A)\) 表示有限集 \(A\) 中元素的个数.

\begin{circlelist}
\item 对任意有限集 \(A\)，\(B\)，“\(A \neq B\)”是“\(d(A, B) > 0\)”的充分必要条件；

\item 对任意有限集 \(A\)，\(B\)，\(C\)，\(d(A, C) \le d(A, B) + d(B, C)\).
\end{circlelist}

（\quad）

\choices
    {命题\circled{1}和命题\circled{2}都成立}
    {命题\circled{1}和命题\circled{2}都不成立}
    {命题\circled{1}成立，命题\circled{2}不成立}
    {命题\circled{1}不成立，命题\circled{2}成立}
\end{problem}

\begin{answer}
A
\end{answer}

\begin{solution}
由集合的并、交关系，\(d(A,B)\) 等于 \(A\) 中不属于 \(B\) 的元素个数与 \(B\) 中不属于 \(A\) 的元素个数之和，即 \(d(A,B)\) 是对称差 \(A\mathbin{\triangle}B\) 的元素个数. 因此 \(d(A,B)=0\) 当且仅当 \(A\subseteq B\) 且 \(B\subseteq A\)，也就是 \(A=B\)，命题\(\circled{1}\)成立.

又 \(A\mathbin{\triangle}C\subseteq(A\mathbin{\triangle}B)\cup(B\mathbin{\triangle}C)\)，所以
\(d(A,C)\leqslant d(A,B)+d(B,C)\)，命题\circled{2}也成立. 因此选 A.
\end{solution}

\begin{problem}
存在函数 \( f(x) \) 满足：对于任意 \( x \in \R \) 都有（\quad）

\choices
    {\(f(\sin 2x) = \sin x\)}
    {\(f(\sin 2x) = x^2 + x\)}
    {\(f(x^{2} + 1) = |x + 1|\)}
    {\(f(x^{2} + 2x) = |x + 1|\)}
\end{problem}

\begin{answer}
D
\end{answer}

\begin{solution}
若 \(f(\sin2x)=\sin x\) 对任意 \(x\) 成立，取 \(x=\dfrac{\pi}{6}\) 和 \(x=\dfrac{\pi}{3}\)，两者的 \(\sin2x\) 相同而 \(\sin x\) 不同，矛盾.

若 \(f(\sin2x)=x^2+x\) 对任意 \(x\) 成立，取 \(x=0\) 和 \(x=\dfrac{\pi}{2}\)，两者的 \(\sin2x\) 都为 \(0\)，而 \(x^2+x\) 的值不同，矛盾.

若 \(f(x^2+1)=|x+1|\) 对任意 \(x\) 成立，取 \(x=1\) 和 \(x=-1\)，两者的 \(x^2+1\) 相同而 \(|x+1|\) 不同，矛盾.

对于选项 D，有 \(x^2+2x=(x+1)^2-1\)，且 \(|x+1|=\sqrt{x^2+2x+1}\). 例如定义
\(f(t)=\begin{cases}
\sqrt{t+1},&t\geqslant-1,\\
0,&t<-1,
\end{cases}\)
则对任意 \(x\in\R\)，\(f(x^2+2x)=|x+1|\). 因此选 D.
\end{solution}

\begin{problem}
如图，已知 \(\triangle ABC\)，\(D\) 是 \(AB\) 的中点，沿直线 \(CD\) 将 \(\triangle ACD\) 翻折成 \(\triangle A'CD\)，所成二面角 \(A' - CD - B\) 的平面角为 \(\alpha\)，则（\quad）

\bitmapfigure[max width=0.52\linewidth]{img/2015/zhejiang_science/q8_fig1.png}

\choices
    {\(\angle A'DB \leqslant \alpha\)}
    {\(\angle A'DB \geqslant \alpha\)}
    {\(\angle A'CB \leqslant \alpha\)}
    {\(\angle A'CB \geqslant \alpha\)}
\end{problem}

\begin{answer}
B
\end{answer}

\begin{solution}
设 \(\theta=\angle A'DC\). 折叠不改变 \(\triangle ACD\) 的形状，所以 \(\angle A'DC=\angle ADC=\theta\). 又 \(A\)，\(D\)，\(B\) 共线，故 \(\angle BDC=\pi-\theta\).

在过 \(D\) 且垂直于 \(CD\) 的平面内，分别取 \(DA'\)，\(DB\) 在该平面上的投影. 根据二面角平面角的定义，这两个投影的夹角为 \(\alpha\). 将 \(DA'\)，\(DB\) 的单位方向向量沿 \(CD\) 和垂直于 \(CD\) 的方向分解，得
\(\cos\angle A'DB=-\cos^2\theta+\sin^2\theta\cos\alpha\).

因此
\(\cos\angle A'DB-\cos\alpha=-\cos^2\theta(1+\cos\alpha)\leqslant0\). 由于角均取 \(0\) 到 \(\pi\) 之间的值，得到 \(\angle A'DB\geqslant\alpha\). 因此选 B.
\end{solution}

\section{填空题}

\begin{problem}
双曲线 \(\frac{x^2}{2} - y^2 = 1\) 的焦距是\fillinblank{}，渐近线方程是\fillinblank{}.
\end{problem}

\begin{answer}
\(2\sqrt{3}\)，\(y=\pm\dfrac{\sqrt{2}}{2}x\)
\end{answer}

\begin{solution}
将方程写成标准形式
\(\dfrac{x^2}{(\sqrt{2})^2}-\dfrac{y^2}{1^2}=1\)，可知半实轴 \(a=\sqrt{2}\)，半虚轴 \(b=1\). 因此
\(c=\sqrt{a^2+b^2}=\sqrt{3}\)，焦距为 \(2c=2\sqrt{3}\).

双曲线的渐近线方程为 \(y=\pm\dfrac{b}{a}x\)，所以渐近线为
\(y=\pm\dfrac{\sqrt{2}}{2}x\).
\end{solution}

\begin{problem}
已知函数 \( f(x) = \begin{cases}
x + \frac{2}{x} - 3, & x \geqslant 1,\\
\lg (x^2 + 1), & x < 1,
\end{cases} \) 则 \( f(f(-3)) =\fillinblank{} \)，\( f(x) \) 的最小值是\fillinblank{}.
\end{problem}

\begin{answer}
\(0\)，\(2\sqrt{2}-3\)
\end{answer}

\begin{solution}
因为 \(-3<1\)，所以 \(f(-3)=\lg(10)=1\)，从而
\(f(f(-3))=f(1)=1+\dfrac{2}{1}-3=0\).

当 \(x<1\) 时，\(x^2+1\geqslant1\)，故 \(\lg(x^2+1)\geqslant0\)，且在 \(x=0\) 时取到 \(0\). 当 \(x\geqslant1\) 时，
\(x+\dfrac{2}{x}-3=\left(\sqrt{x}-\sqrt{\dfrac{2}{x}}\right)^2+2\sqrt{2}-3\geqslant2\sqrt{2}-3\)，等号在 \(x=\sqrt{2}\) 时成立. 又 \(2\sqrt{2}-3<0\)，故 \(f(x)\) 的最小值为 \(2\sqrt{2}-3\).
\end{solution}

\begin{problem}
函数 \( f(x) = \sin^2 x + \sin x \cos x + 1 \) 的最小正周期是\fillinblank{}，单调递减区间是\fillinblank{}.
\end{problem}

\begin{answer}
\(\pi\)，\(\left[\frac{3\pi}{8}+k\pi,\frac{7\pi}{8}+k\pi\right]\)（\(k\in\Z\)）
\end{answer}

\begin{solution}
利用二倍角公式，
\(f(x)=\dfrac{3}{2}+\dfrac{1}{2}(\sin2x-\cos2x)=\dfrac{3}{2}+\dfrac{\sqrt{2}}{2}\sin\left(2x-\dfrac{\pi}{4}\right)\). 因此 \(f(x)\) 的最小正周期为 \(\pi\).

又
\(f'(x)=\sin2x+\cos2x=\sqrt{2}\sin\left(2x+\dfrac{\pi}{4}\right)\). 当
\(\dfrac{3\pi}{8}+k\pi<x<\dfrac{7\pi}{8}+k\pi\) 时，\(f'(x)<0\)，故单调递减区间可写为
\(\left[\dfrac{3\pi}{8}+k\pi,\dfrac{7\pi}{8}+k\pi\right]\)，其中 \(k\in\Z\).
\end{solution}

\begin{problem}
若 \(a = \log_{4} 3\)，则 \(2^a + 2^{-a} =\fillinblank{}.\)
\end{problem}

\begin{answer}
\(\frac{4\sqrt{3}}{3}\)
\end{answer}

\begin{solution}
由 \(a=\log_4 3\)，得 \(4^a=3\)，即 \(2^{2a}=3\). 因为 \(2^a>0\)，所以 \(2^a=\sqrt{3}\)，从而
\(2^a+2^{-a}=\sqrt{3}+\dfrac{1}{\sqrt{3}}=\dfrac{4\sqrt{3}}{3}\).
\end{solution}

\begin{problem}
如图，在三棱锥 \(A - BCD\) 中，\(AB = AC = BD = CD = 3\)，\(AD = BC = 2\)，点 \(M\)，\(N\) 分别为 \(AD\)，\(BC\) 的中点，则异面直线 \(AN\)，\(CM\) 所成的角的余弦值是\fillinblank{}.

\bitmapfigure[max width=0.38\linewidth]{img/2015/zhejiang_science/q13_fig1.png}
\end{problem}

\begin{answer}
\(\frac{7}{8}\)
\end{answer}

\begin{solution}
取 \(A\) 为原点，设 \(\bs{p}=\overrightarrow{AB}\)，\(\bs{q}=\overrightarrow{AC}\)，\(\bs{r}=\overrightarrow{AD}\). 由已知边长和数量积公式，
\(\bs{p}\cdot\bs{q}=\dfrac{3^2+3^2-2^2}{2}=7\)，
\(\bs{p}\cdot\bs{r}=\dfrac{3^2+2^2-3^2}{2}=2\)，
\(\bs{q}\cdot\bs{r}=\dfrac{3^2+2^2-3^2}{2}=2\).

由于 \(N\)，\(M\) 分别为 \(BC\)，\(AD\) 的中点，
\(\overrightarrow{AN}=\dfrac{\bs{p}+\bs{q}}{2}\)，
\(\overrightarrow{CM}=\dfrac{\bs{r}-2\bs{q}}{2}\). 因而
\(\overrightarrow{AN}\cdot\overrightarrow{CM}
=\dfrac{1}{4}(\bs{p}+\bs{q})\cdot(\bs{r}-2\bs{q})
=\dfrac{1}{4}(2-14+2-18)=-7\).

同时
\(|\overrightarrow{AN}|^2=\dfrac{1}{4}|\bs{p}+\bs{q}|^2=\dfrac{1}{4}(9+9+14)=8\)，
\(|\overrightarrow{CM}|^2=\dfrac{1}{4}|\bs{r}-2\bs{q}|^2=\dfrac{1}{4}(4+36-8)=8\).
异面直线所成的角取锐角，故其余弦值为
\(\dfrac{|\overrightarrow{AN}\cdot\overrightarrow{CM}|}{|\overrightarrow{AN}|\,|\overrightarrow{CM}|}=\dfrac{7}{8}\).
\end{solution}

\begin{problem}
若实数 \(x\)，\(y\) 满足 \(x^{2} + y^{2} \leqslant 1\)，则 \(\left|2x + y - 2\right| + \left|6 - x - 3y\right|\) 的最小值是\fillinblank{}.
\end{problem}

\begin{answer}
\(3\)
\end{answer}

\begin{solution}
令 \(u=2x+y-2\)，\(v=6-x-3y\). 由绝对值不等式，
\(|u|+|v|\geqslant-u+v=8-3x-4y\).

又由柯西不等式及 \(x^2+y^2\leqslant1\)，
\(3x+4y\leqslant\sqrt{3^2+4^2}\sqrt{x^2+y^2}\leqslant5\). 因此
\(|2x+y-2|+|6-x-3y|\geqslant8-5=3\).

取 \(x=\dfrac{3}{5}\)，\(y=\dfrac{4}{5}\)，此时 \(x^2+y^2=1\)，
\(2x+y-2=0\)，\(6-x-3y=3\)，原式等于 \(3\). 所以最小值为 \(3\).
\end{solution}

\begin{problem}
已知 \(\bs{e}_{1}\)，\(\bs{e}_{2}\) 是空间单位向量，\(\bs{e}_{1} \cdot \bs{e}_{2} = \frac{1}{2}\)，若空间向量 \( \bs{b} \) 满足 \( \bs{b} \cdot \bs{e}_{1} = 2\)，\( \bs{b} \cdot \bs{e}_{2} = \frac{5}{2}\)，且对于任意 \(x\)，\(y \in \R\)，\(\left|\bs{b} - (x \bs{e}_{1} + y \bs{e}_{2})\right| \geqslant \left|\bs{b} - (x_{0} \bs{e}_{1} + y_{0} \bs{e}_{2})\right| = 1\)（\(x_{0}, y_{0} \in \R\)），则 \(x_{0} =\)\fillinblank{}，\(y_{0} =\)\fillinblank{}，\(|\bs{b}| =\)\fillinblank{}.
\end{problem}

\begin{answer}
\(1\)，\(2\)，\(2\sqrt{2}\)
\end{answer}

\begin{solution}
设
\(\bs{r}=\bs{b}-(x_0\bs{e}_1+y_0\bs{e}_2)\). 由题意，\(\bs{r}\) 是从 \(\bs{b}\) 到平面
\(\{x\bs{e}_1+y\bs{e}_2\mid x,y\in\R\}\) 的最短向量，所以
\(\bs{r}\perp\bs{e}_1\) 且 \(\bs{r}\perp\bs{e}_2\). 因此
\(0=\bs{r}\cdot\bs{e}_1=2-x_0-\dfrac{y_0}{2}\)，
\(0=\bs{r}\cdot\bs{e}_2=\dfrac{5}{2}-\dfrac{x_0}{2}-y_0\).

解方程组
\(\begin{cases}
x_0+\dfrac{y_0}{2}=2,\\
\dfrac{x_0}{2}+y_0=\dfrac{5}{2}
\end{cases}\)，得 \(x_0=1\)，\(y_0=2\).

又 \(|\bs{r}|=1\)，且 \(\bs{r}\perp\bs{e}_1+2\bs{e}_2\)，所以
\(|\bs{b}|^2=|\bs{e}_1+2\bs{e}_2|^2+|\bs{r}|^2
=1+4+4\cdot\dfrac{1}{2}+1=8\). 因而 \(|\bs{b}|=2\sqrt{2}\).
\end{solution}

\section{解答题}

\begin{problem}
在 \(\triangle ABC\) 中，内角 \(A\)，\(B\)，\(C\) 所对的边分别是 \(a\)，\(b\)，\(c\). 已知 \(A = \frac{\pi}{4}\)，\(b^2 - a^2 = \frac{1}{2} c^2\).
\begin{enumerate}
\item 求 \(\tan C\) 的值；
\item 若 \(\triangle ABC\) 的面积为 3，求 \(b\) 的值.
\end{enumerate}
\end{problem}

\begin{solution}
\begin{enumerate}
\item 由余弦定理，\(a^2=b^2+c^2-2bc\cos A=b^2+c^2-\sqrt{2}bc\)，所以
\(b^2-a^2=\sqrt{2}bc-c^2\). 结合已知条件得
\(\sqrt{2}bc-c^2=\frac{1}{2}c^2\)，从而 \(b=\frac{3}{2\sqrt{2}}c\).

又由正弦定理，\(\dfrac{\sin B}{\sin C}=\dfrac{b}{c}=\dfrac{3}{2\sqrt{2}}\)，且 \(B=\dfrac{3\pi}{4}-C\)，故
\(\dfrac{\sqrt{2}}{2}(\cos C+\sin C)=\dfrac{3}{2\sqrt{2}}\sin C\). 因而 \(2\cos C=\sin C\)，所以 \(\tan C=2\).

\item 由三角形面积公式，\(3=\dfrac{1}{2}bc\sin A=\dfrac{\sqrt{2}}{4}bc\). 代入 \(b=\dfrac{3}{2\sqrt{2}}c\)，得 \(3=\dfrac{3}{8}c^2\)，所以 \(c=2\sqrt{2}\)，从而 \(b=3\).
\end{enumerate}
\end{solution}

\begin{problem}
如图，在三棱柱 \(ABC - A_{1}B_{1}C_{1}\) 中，\(\angle BAC = 90^{\circ}\)，\(AB = AC = 2\)，\(A_{1}A = 4\)，\(A_{1}\) 在底面 \(ABC\) 的射影为 \(BC\) 的中点，\(D\) 是 \(B_{1}C_{1}\) 的中点.
\begin{enumerate}
\item 证明： \(A_{1}D \perp\) 平面 \(A_{1}BC\)；
\item 求二面角 \(A_{1} - BD - B_{1}\) 的平面角的余弦值.
\end{enumerate}

\bitmapfigure[max width=0.42\linewidth]{img/2015/zhejiang_science/q17_fig1.png}
\end{problem}

\begin{solution}
\begin{enumerate}
\item 设 \(O\) 为 \(BC\) 的中点. 由 \(AB=AC=2\)，\(\angle BAC=90^{\circ}\)，得 \(BC=2\sqrt{2}\)，\(AO=\sqrt{2}\). 由于 \(A_1\) 在底面 \(ABC\) 的射影为 \(O\)，所以 \(A_1O\perp\) 平面 \(ABC\)，且在直角三角形 \(AA_1O\) 中，\(A_1A=4\)，从而 \(A_1O=\sqrt{14}\).

以 \(OB\)，\(OA\)，\(OA_1\) 所在直线分别为 \(x\)，\(y\)，\(z\) 轴建立空间直角坐标系，则
\(A_1(0,0,\sqrt{14})\)，\(B(\sqrt{2},0,0)\)，\(C(-\sqrt{2},0,0)\). 由三棱柱的平移关系，\(D\) 为 \(B_1C_1\) 的中点，故 \(D(0,-\sqrt{2},\sqrt{14})\).

于是 \(\overrightarrow{A_1D}=(0,-\sqrt{2},0)\)，而平面 \(A_1BC\) 内的两条相交直线可取 \(BC\) 和 \(OA_1\)，它们的方向向量分别为 \((-2\sqrt{2},0,0)\) 和 \((0,0,\sqrt{14})\). 这两个向量都与 \(\overrightarrow{A_1D}\) 垂直，所以 \(A_1D\perp\) 平面 \(A_1BC\).

\item 取 \(\bs{u}=\overrightarrow{BD}=(-\sqrt{2},-\sqrt{2},\sqrt{14})\)，则 \(|\bs{u}|^2=18\). 将 \(\overrightarrow{BA_1}=(-\sqrt{2},0,\sqrt{14})\) 和 \(\overrightarrow{BB_1}=(0,-\sqrt{2},\sqrt{14})\) 分别投影到垂直于 \(BD\) 的方向上，得到平面 \(A_1BD\) 和平面 \(B_1BD\) 内分别指向 \(A_1\) 和 \(B_1\) 的向量
\[
\bs{p}=\overrightarrow{BA_1}-\frac{\overrightarrow{BA_1}\cdot\bs{u}}{|\bs{u}|^2}\bs{u}=\left(-\frac{\sqrt{2}}{9},\frac{8\sqrt{2}}{9},\frac{\sqrt{14}}{9}\right),
\quad
\bs{q}=\overrightarrow{BB_1}-\frac{\overrightarrow{BB_1}\cdot\bs{u}}{|\bs{u}|^2}\bs{u}=\left(\frac{8\sqrt{2}}{9},-\frac{\sqrt{2}}{9},\frac{\sqrt{14}}{9}\right),
\]
其中 \(\overrightarrow{BA_1}\cdot\bs{u}=\overrightarrow{BB_1}\cdot\bs{u}=16\). 直接计算得 \(\bs{p}\cdot\bs{q}=-\frac{2}{9}\)，且 \(|\bs{p}|=|\bs{q}|=\frac{4}{3}\). 因而二面角 \(A_1-BD-B_1\) 的平面角余弦为
\(\dfrac{\bs{p}\cdot\bs{q}}{|\bs{p}|\,|\bs{q}|}=-\dfrac{1}{8}\).
\end{enumerate}
\end{solution}

\begin{problem}
已知函数 \( f(x) = x^2 + ax + b \) （\(a\)，\(b \in \R\)），记 \( M(a, b) \) 是 \( |f(x)| \) 在区间 \([-1, 1]\) 上的最大值.
\begin{enumerate}
\item 证明：当 \(|a| \geqslant 2\) 时，\(M(a, b) \geqslant 2\)；
\item 当 \(a\)，\(b\) 满足 \(M(a, b) \leqslant 2\) 时，求 \(|a| + |b|\) 的最大值.
\end{enumerate}
\end{problem}

\begin{solution}
\begin{enumerate}
\item 因为 \(f(1)-f(-1)=2a\)，所以
\(\max\{|f(1)|,|f(-1)|\}\geqslant\dfrac{|f(1)-f(-1)|}{2}=|a|\). 当 \(|a|\geqslant2\) 时，\(M(a,b)\geqslant\max\{|f(1)|,|f(-1)|\}\geqslant2\).

\item 由 \(M(a,b)\leqslant2\)，得 \(|b|=|f(0)|\leqslant2\)，且
\[
|a|+|1+b|=\max\{|1+b+a|,|1+b-a|\}\leqslant2.
\]
当 \(b\geqslant0\) 时，\(|a|+|b|\leqslant1\). 当 \(-1\leqslant b<0\) 时，由 \(|a|\leqslant1-b\)，得 \(|a|+|b|\leqslant1-2b\leqslant3\). 当 \(-2\leqslant b<-1\) 时，由 \(|a|\leqslant3+b\)，得 \(|a|+|b|\leqslant3\). 因此总有 \(|a|+|b|\leqslant3\).

取 \(a=2\)，\(b=-1\)，则 \(f(x)=x^2+2x-1=(x+1)^2-2\). 当 \(-1\leqslant x\leqslant1\) 时，\(0\leqslant(x+1)^2\leqslant4\)，故 \(|f(x)|\leqslant2\)，即 \(M(2,-1)=2\)，且 \(|a|+|b|=3\). 所以最大值为 \(3\).
\end{enumerate}
\end{solution}

\begin{problem}
已知椭圆 \(\frac{x^2}{2} + y^2 = 1\) 上两个不同的点 \(A\)，\(B\) 关于直线 \(y = mx + \frac{1}{2}\) 对称.
\begin{enumerate}
\item 求实数 \(m\) 的取值范围.
\item 求 \( \triangle AOB \) 面积的最大值（\(O\) 为坐标原点）.
\end{enumerate}
\end{problem}

\begin{solution}
\begin{enumerate}
\item 设 \(P=(u,mu+\frac{1}{2})\) 为 \(A\)，\(B\) 的中点. 由于对称轴垂直于线段 \(AB\)，存在 \(t>0\)，使得
\(A=P+t(m,-1)\)，\(B=P-t(m,-1)\). 将这两点代入椭圆方程并相减，得到 \(um=2y\)，其中 \(y=mu+\frac{1}{2}\). 因而 \(mu=-1\)，所以 \(m\neq0\)，且
\(P=\left(-\frac{1}{m},-\frac{1}{2}\right)\).

记 \(Q(x,y)=\frac{x^2}{2}+y^2\)，记 \(v=(m,-1)\). 由中点条件，二次项的交叉项为零，故
\[
t^2=\frac{1-Q(P)}{Q(v)}=\frac{1-\frac{1}{2m^2}-\frac14}{\frac{m^2}{2}+1}=\frac{3m^2-2}{2m^2(m^2+2)}.
\]
因为 \(t>0\)，必须且只需 \(m^2>\frac{2}{3}\)，所以 \(m\) 的取值范围为 \(\left(-\infty,-\frac{\sqrt{6}}{3}\right)\cup\left(\frac{\sqrt{6}}{3},+\infty\right)\).

\item 设 \(T\) 为 \(\triangle AOB\) 的面积，则
\(T=\frac12|\det(A,B)|=t|\det(P,v)|=t\frac{m^2+2}{2|m|}\). 令 \(z=m^2\)，由上式得
\[
T^2=\frac{(3z-2)(z+2)}{8z^2},\qquad z>\frac23.
\]
又 \(T^2-\frac12=-\dfrac{(z-2)^2}{8z^2}\leqslant0\)，当且仅当 \(z=2\)，即 \(m=\pm\sqrt{2}\) 时取等号. 因此 \(\triangle AOB\) 面积的最大值为 \(\frac{\sqrt{2}}{2}\).
\end{enumerate}
\end{solution}

\begin{problem}
已知数列 \(\{a_{n}\}\) 满足 \(a_1 = \frac{1}{2}\) 且 \(a_{n + 1} = a_n - a_n^2\)（\(n\in \N^*\)）.
\begin{enumerate}
\item 证明： \(1 \leqslant \frac{a_n}{a_{n+1}} \leqslant 2\)（\(n \in \N^*\)）；
\item 设数列 \(\{a_n^2\}\) 的前 \(n\) 项和为 \(S_n\)，证明： \(\frac{1}{2(n+2)} \leqslant \frac{S_n}{n} \leqslant \frac{1}{2(n+1)}\) （\(n \in \N^*\)）.
\end{enumerate}
\end{problem}

\begin{solution}
\begin{enumerate}
\item 由递推式得 \(a_{n+1}=a_n(1-a_n)\). 先证 \(0<a_n\leqslant\frac12\). 当 \(n=1\) 时结论成立；若 \(0<a_n\leqslant\frac12\)，则 \(0<a_{n+1}=a_n(1-a_n)\leqslant\frac14\leqslant\frac12\)，所以对任意 \(n\in\N^*\) 都有该范围. 因而
\(\dfrac{a_n}{a_{n+1}}=\dfrac{1}{1-a_n}\)，且 \(\frac12\leqslant1-a_n<1\)，从而 \(1\leqslant\dfrac{a_n}{a_{n+1}}\leqslant2\).

\item 令 \(g(x)=x(1-x)\). 在 \(\left[0,\frac12\right]\) 上，\(g'(x)=1-2x\geqslant0\)，所以 \(g\) 递增. 证明对任意 \(n\in\N^*\) 有
\[
\frac{1}{2n}\leqslant a_n\leqslant\frac{1}{n+1}.
\]
当 \(n=1\) 时成立. 若对某个 \(n\) 成立，则由 \(g\) 的单调性，
\[
a_{n+1}=g(a_n)\geqslant g\left(\frac{1}{2n}\right)=\frac{2n-1}{4n^2}\geqslant\frac{1}{2(n+1)},
\quad
a_{n+1}=g(a_n)\leqslant g\left(\frac{1}{n+1}\right)=\frac{n}{(n+1)^2}\leqslant\frac{1}{n+2}.
\]
故上述不等式对所有正整数 \(n\) 成立. 又由 \(a_k-a_{k+1}=a_k^2\)，数列 \(\{a_n^2\}\) 的前 \(n\) 项和满足
\(S_n=\sum_{k=1}^{n}a_k^2=\sum_{k=1}^{n}(a_k-a_{k+1})=\frac12-a_{n+1}\). 因此
\[
\frac{n}{2(n+2)}=\frac12-\frac{1}{n+2}\leqslant S_n\leqslant\frac12-\frac{1}{2(n+1)}=\frac{n}{2(n+1)}.
\]
两边同除以正数 \(n\)，即得 \(\dfrac{1}{2(n+2)}\leqslant\dfrac{S_n}{n}\leqslant\dfrac{1}{2(n+1)}\).
\end{enumerate}
\end{solution}
